\documentclass{article}
\usepackage{/globals/de}
\usepackage{/globals/anki}
\ankiprefix{wissensmanagement}
\ankitopic{Wissensmanagement}
\ankishow
\usepackage{csquotes}
\begin{document}
\section{Wissensmanagement} 
Es ist wichtig dass kein sog. \textquote{Kopfmonopol} aufgebaut wird; dass es dazu kommen kann dass bestimmte Prozesse durchgeführt werden können, einzig und allein weil die eine Person mit dem notwendigen Wissen fehlt. Hierzu dient gutes Wissensmanagement.
\ankinote{kopfmonopol}{Kopfmonopol}{Eine Person ohne dessen Wissen Prozesse nicht funktionieren}

\subsection{Wissensarten}
Wissen besteht in zwei Arten
\begin{description}
 \item[Explizites Wissen] Wissen welches bereits aufgeschrieben wurde. Es ist personenunabhängig.
 \item[Implizites Wissen] Das Wissen das nirgendswo aufgeschrieben ist und auch nur schwer aufzuschreieben ist, sondern nur in den Köpfen mancher ArbeiterInnen aufzufinden ist. Für diese Personen ist es oftmals schneller aufrufbar.
\end{description}
\ankinote{wissensarten}{Wissensarten}{Explizit Wissen, Implizites Wissen}
\ankinote{wissensarten-explizit}{Wissensarten Explizit}{Aufgeschriebenes Wissen}
\ankinote{wissensarten-implizit}{Wissensarten Implizit}{Nur im Kopf, schwer in Worte zu fassen, Intuition}

\subsection{Wissenskreislauf}
\begin{description}
 \item[Sozialisation] Durch gemeinsame Arbeite und Gespräche entsteht implizites Wissen. Hierbei können auch Hospitationen oder Jobrotationen genutzt werden
 \item[Externalisierung] Implizites Wissen wird gezielt ausgeschrieben. Dies kann in Dokumenten, Mindmaps, Wikis, Videos, \etc passieren.
 \item[Kombination] Explizites Wissen wird zusammengeführt und strukturiert.
 \item[Internalisierung] Explizites Wissen wird angewendet, bis es implizit ist.
\end{description}
\ankinote{wissenskreislauf}{Wissenskreilauf}{}
\ankinote{wissenskreislauf-soz}{Wissenskreilauf Sozialisation}{Gemeinsame Arbeit und Gespräche, Hospitationen, Jobrotationen entsteht Wissen}
\ankinote{wissenskreislauf-ext}{Wissenskreilauf Externalisierung}{Wissen aufschreiben}
\ankinote{wissenskreislauf-kom}{Wissenskreilauf Kombination}{Zusammenführen, strukturieren}
\ankinote{wissenskreislauf-int}{Wissenskreilauf Internalisierung}{Anwenden, bis es implizit ist}

\subsection{Herausfoderungen}
Mancher ArbeiterInnen wollen ihr Wissen nicht wirklich Teilen, weil dies sie ersetzbar machen würde. Zudem ist die Zeit die mit dem Wissensmanagement verbracht wird auch Zeit die nicht mit direkt wirtschaftlich profitabler Arbeit verbracht wird.
\ankinote{herausforderungen}{Herausforderungen}{Hemmungen, ersetzbarkeit; Arbeitszeit}
\end{document}